\documentclass[11pt, oneside]{article}
\usepackage{geometry}
\geometry{letterpaper}
\usepackage{graphicx}	
\usepackage{amssymb}
\usepackage{amsmath}
\usepackage{setspace}
\doublespacing
\setlength{\parindent}{0pt}
 
\begin{document}

\section*{Definition species}

The metabolic parser will generate a window with a first void line (like the existing one) in which the species are defined.

Species are defined at the beginning, like $A, B, C, ...$

At the end of this line there is a CONFIRM button that fixes the species. If the user forgets one species, it is enough to add it also later, and click again "confirm".

After "confirm", a first empty reaction (a line) is created. Each reaction contains 4 FIELDS and 2 buttons.

FIELD1 is a drop-down menu presenting a choice among all the species, it is the "source".

A similar FIELD3 with the same drop-down menu defines the "target" of the relation.

 \section*{Definition actions}

Each line defines an action between species in the form of a $verbal$. The list of possible verbals appears in a drop-down menu FIELD 2, placed between the two above.

The verbals in the drop-down FIELD2 are: PRODUCES, CONSUMES, ACTIVATES.

 \section*{Definition kinetics}
 
 \begin{itemize}
\item The detailed kinetics of the relation is specified in FIELD4, another drop-down menu including the list of kinetic keywords in Section 2 below: LINEAR, MASS ACTION, ACTIVATION (MM), ACTIVATION (HILL), INHIBITION (HILL), INHIBITION (COMP), REVERSIBLE.

\item Some of them include numerical parameter(s) that the user is asked to specify before moving to the next reaction.

\item Some kinetics require a third partner, that appears in a FIELD5 as drop-down menu with the species. This is optional and can be left blank if the kinetics in FIELD4 does not require it.

\item Each of these equations can be modified by multiplying by other terms, in a chain. This is done with the GATE button.

The 6th field is the GATE button. If clicked, it will open a second line below the current one. It takes as input (in FIELD3) the previous equation, instead of one species. A species in FIELD1 can be selected, and will perform an action/kinetics that will multiply (activation/repression) the equation now appearing as FIELD3. This process can be repeated to describe a contribution depending on the product of 2, 3, 4... activation/repression multipliers.

\item The 7th field is the COMMIT button. When clicked this indicates the user is happy to move on to the next reaction. 

\item (Optional : Possibly, a printout below the first line could summarize the result in real time as it is modified.) 

\end{itemize} 

So, each new line looks like:

FIELD1 (menu) FIELD2 (menu) FIELD3 (menu) FIELD4 (menu) FIELD5 (menu opt) GATE (button) COMMIT (button)

\begin{itemize}
\item Once all the reactions are specified in the N lines, the result will be a number of different "flux terms" contributing to each species.

\item *** BE CAREFUL HERE *** in case the flux term is the result of a product $v \cdot f1 \cdot f2 ...$ with successive applications of GATE, to count in the contribution only the final result, and not sum each term individually.

\item A final button like CONFIRM concludes the process and tells the interpreter to generate the corresponding equations $dA/dt, dB/dt, dC/dt, ...$ by regrouping all the contributing flux terms from the different reaction lines.
\end{itemize}

\section*{Minimal catalog of kinetic keywords}

\subsection*{Definition of allowed actions/kinetics}

In the following FIELD1 = A, FIELD3 = B, FIELD5 =C (if needed).

\begin{itemize}
\item Verbal actions (FIELD2) are paired with kinetic choices (FIELD4). 
\item Only some combinations FIELD2/FIELD4 are allowed. 
\item With the GATE option, a pair verbal/kinetics can be multiplied in cascade by other pairs verbal/kinetics: the difference is that in the first definition FIELD1 and FIELD3 are both species; in the GATE definition FIELD1 is still a species and FIELD3 is fixed as the previous reaction.
\end{itemize}


\subsection*{Explicit equations}

Below there is the pairing between fields and the flux terms they generate.

\subsection{FIELD2 = PRODUCES}

\subsubsection{FIELD4 = LINEAR}

Flux term $v = kA$ 

Contributions
\begin{align}
\frac{dA}{dt} = -v \\
\frac{dB}{dt} = +v 
\end{align}

Present the prompt: "$v = kA$, Define a value for $k$". Read $k$ from an input field.
Only after field is valorized can click GATE or COMMIT.


\subsubsection{FIELD4 = MASS ACTION}

Flux term $v = kAB$ 

Contributions
\begin{align}
\frac{dA}{dt} = -v \\
\frac{dB}{dt} = -v \\
\frac{dC}{dt} = +v
\end{align}

Present the prompt: "$v = kAB$, Define a value for $k$". Read $k$ from an input field.
Only after field is valorized can click GATE or COMMIT.


\subsubsection{FIELD4 = ACTIVATION (MM)}

Flux term $v = V \tfrac{ A }{ K + A}$ 

Contributions
\begin{align}
\frac{dA}{dt} = -v \\
\frac{dB}{dt} = +v
\end{align}

Present the prompt: "$v = V A/(K+A)$, Define a value for $V,K$". Read $V,K$ from two input fields.
Only after two fields are valorized can click GATE or COMMIT.


\subsubsection{FIELD4 = ACTIVATION (HILL)}

Flux term $v = V \tfrac{ A^n }{ K^n + A^n}$,

\begin{align}
\frac{dA}{dt} = -v \\
\frac{dB}{dt} = +v
\end{align}

Present the prompt: "$v = V A^n/(K^n+A^n)$, Define a value for $V,K,n$". Read $V,K,n$ from three input fields.
Only after the three fields are valorized can click GATE or COMMIT.


\subsubsection{FIELD4 = INHIBITION (HILL)}

Flux multiplier $f = V \tfrac{ K^n }{ K^n + A^n}$

It can only appear as multiplier of an already defined flux term $v$ (with GATE), it has no independent contribution:
\begin{align}
\frac{dA}{dt} = v \cdot f \\
\frac{dB}{dt} = v \cdot f
\end{align}

Present the prompt: "$f = V A^n/(K^n+A^n)$, Define a value for $V,K,n$". Read $V,K,n$ from three input fields.
Only after the three fields are valorized can click GATE or COMMIT.


\subsubsection{FIELD4 = REVERSIBLE}

Flux term $v = k_1 A - k_2 B$,

\begin{align}
\frac{dA}{dt} = -v \\
\frac{dB}{dt} = +v
\end{align}

Present the prompt: "$v = k_1 A - k_2 B$, Define a value for $k1,k2$". Read $k1,k2$ from two input fields.
Only after the two fields are valorized can click GATE or COMMIT.

\subsection{FIELD2 = CONSUMES}

\subsubsection{FIELD4 = LINEAR}

Flux term $v = kAB$ 

Contributions
\begin{align}
\frac{dA}{dt} = 0 \\
\frac{dB}{dt} = -v 
\end{align}

Present the prompt: "$v = kB$, Define a value for $k$". Read $k$ from an input field.
Only after field is valorized can click GATE or COMMIT.

\subsubsection{FIELD4 = MASS ACTION}

Flux term $v = kAB$ 

Contributions
\begin{align}
\frac{dA}{dt} = -v \\
\frac{dB}{dt} = -v 
\end{align}

Present the prompt: "$v = kB$, Define a value for $k$". Read $k$ from an input field.
Only after field is valorized can click GATE or COMMIT.


\subsubsection{FIELD4 = ACTIVATION (MM)}

Flux term $v = V \tfrac{ B }{ K + B}$ 

Contributions
\begin{align}
\frac{dA}{dt} = 0 \\
\frac{dB}{dt} = -v
\end{align}

Present the prompt: "$v = V B/(K+B)$, Define a value for $V,K$". Read $V,K$ from two input fields.
Only after two fields are valorized can click GATE or COMMIT.


\subsection{FIELD2 = ACTIVATES}

\subsubsection{FIELD4 = LINEAR}

Flux term $v = kA$ 

Contributions
\begin{align}
\frac{dA}{dt} = 0 \\
\frac{dB}{dt} = +v 
\end{align}

Present the prompt: "$v = kA$, Define a value for $k$". Read $k$ from an input field.
Only after field is valorized can click GATE or COMMIT.


\subsubsection{FIELD4 = ACTIVATION (MM)}

Flux term $v = V \tfrac{ A }{ K + A }$ 

Contributions
\begin{align}
\frac{dA}{dt} = 0 \\
\frac{dB}{dt} = +v
\end{align}

Present the prompt: "$v = V B/(K+B)$, Define a value for $V,K$". Read $V,K$ from two input fields.
Only after two fields are valorized can click GATE or COMMIT.


\subsubsection{FIELD4 = ACTIVATION (HILL)}

Flux term $v = V \tfrac{ A^n }{ K^n + A^n}$,

\begin{align}
\frac{dA}{dt} = 0 \\
\frac{dB}{dt} = +v
\end{align}

Present the prompt: "$v = V A^n/(K^n+A^n)$, Define a value for $V,K,n$". Read $V,K,n$ from three input fields.
Only after the three fields are valorized can click GATE or COMMIT.


\subsubsection{FIELD4 = INHIBITION (HILL)}

Flux multiplier $f = V \tfrac{ K^n }{ K^n + A^n}$

It can only appear as multiplier of an already defined flux term $v$ (with GATE), it has no independent contribution:
\begin{align}
\frac{dA}{dt} = v \cdot f \\
\frac{dB}{dt} = v \cdot f
\end{align}

Present the prompt: "$f = V A^n/(K^n+A^n)$, Define a value for $V,K,n$". Read $V,K,n$ from three input fields.
Only after the three fields are valorized can click GATE or COMMIT.


\subsubsection{FIELD4 = INHIBITION (COMP)}

Flux multiplier $f = V \tfrac{ 1 }{ 1 + aA}$

It can only appear as multiplier of an already defined flux term $v$ (with GATE), it has no independent contribution:
\begin{align}
\frac{dA}{dt} = v \cdot f \\
\frac{dB}{dt} = v \cdot f
\end{align}

Present the prompt: "$f = V/(1+aA)$, Define a value for $V,a$". Read $V,a$ from three input fields.
Only after the two fields are valorized can click GATE or COMMIT.


\end{document}  